% SWE-Bench-CL protocol: the control that makes the comparison mean
% something, the dataset pin, the two disclosed prompt deviations, and
% the retrieval recall that ceilings every arm.
The decisive control is \texttt{random\_matched}: beating
\texttt{memory\_off} alone could be a token-budget effect, so the
learning-curve claim is the second-half-minus-first-half improvement of
\texttt{memory\_on} over \texttt{random\_matched} on at least three
seeds per sequence.

Two properties of this leg bound what it can show, and both are reported
rather than assumed: BM25's recall, below, and the fact that a seed here
does not fix the world it does in the deterministic suites---the endpoint
is sampled, so three seeds are three draws rather than three controlled
replicates (\S\ref{sec:limitations}).

That difference, and not either arm's own second-half-minus-first-half
figure, is the quantity. These sequences are \emph{curricula} ordered by
increasing difficulty, so the second half is harder by construction and
every arm posts a negative within-arm delta whether or not it learned
anything. Both arms walk the identical ordering, so the difficulty
gradient cancels in the paired difference; reading an arm's raw delta as
a learning signal would be reading the curriculum instead.

The dataset is pinned by an exact (commit, sha256) pair; the loader
refuses any file whose hash differs. As of this writing the pin
re-verifies against live upstream, and the two pilot sequences hold
(pytest, 19 tasks; astropy, 22 tasks). The full pipeline is validated
end-to-end \emph{including the real evaluation}, and it resolves real
issues. A minimal prompt (problem statement only) resolves $\approx 0$ for
any model, for two structural reasons: the model never sees the code, and
it cannot compute correct diff hunk line numbers. The run configuration
that produces a curve adds two stdlib pieces, disclosed as a deviation
from the minimal pre-registered prompt: \emph{(i)} BM25 retrieval of the
repository's source files at the task's base commit against the issue text
(no oracle---the gold patch is never read; the query is issue-only and
identical across arms, so arms still differ only in lesson memory), and
\emph{(ii)} a search/replace edit format from which the harness computes
the diff with \texttt{difflib}, so hunk line numbers are correct by
construction. With these, gpt-4.1 resolves \texttt{pytest-dev\_\_pytest-5262}
end-to-end in the official docker harness under \texttt{linux/amd64}
emulation (BM25 retrieved the correct file from the issue text alone; the
edit applied; fail-to-pass 1/1, pass-to-pass 108/108); the same task under
the minimal prompt produced an unappliable diff.

\paragraph{Retrieval recall bounds every arm, so we measured it.}
A file the model never sees is a task no arm can solve, which makes
BM25's recall a ceiling on the whole experiment rather than an
implementation detail. We report it directly: over the pytest sequence,
the fraction of tasks where the file the gold patch edits appears in the
retrieved context is $37\%$ at a 60k-character budget over 5 files,
$63\%$ at 150k/5, $74\%$ at 300k/10, and $89\%$ at 600k/20. The budget
rather than the file count is what binds at the low end---one or two
large modules exhaust it---and the first configuration we ran was the
60k one, where a third of tasks produced no patch at all because the
model wrote edits against files it had not been shown. We report the
matrix at 300k/10 and state the residual: roughly a quarter of tasks
remain unanswerable for retrieval reasons, and they depress every arm
equally. The gold patch is read \emph{only} for this measurement, never
by the retriever and never by any arm; the retrieval query is issue-only
and identical across arms, so recall sets how many tasks are answerable
and never which arm answers them.

